
\paragraph{\textbf{1.12. Proof of \SGARefLink{sga3:xx:1.9}{1.9}.}}\SGAInlineRefTarget{sga3:xx:paragraph:1-12}
We have
\[
  \mathfrak g
    =\mathfrak t\oplus\mathfrak g^\alpha\oplus\mathfrak g^{-\alpha}.
\]
As explained in \SGARefLink{sga3:xx:1.9}{1.9}, regard \(\mathfrak g\) as a
\(T\mathbin{\cdot_\alpha}\mathbb G_{a,k}\)-module. Clearly
\(\mathfrak t\oplus\mathfrak g^\alpha\) is a submodule of
\(\mathfrak g\), and the quotient is isomorphic to
\(\mathfrak g^{-\alpha}\) as a \(k\)-vector space and even as a
\(T\)-module. The group \(\mathbb G_{a,k}\) acts trivially on this
one-dimensional quotient, since every group morphism
\(\mathbb G_{a,k}\to\mathbb G_{m,k}\) is trivial. Likewise,
\(\mathfrak g^\alpha\) is a submodule of
\(\mathfrak t\oplus\mathfrak g^\alpha\); the quotient is isomorphic
to \(\mathfrak t\) as a \(T\)-module, and
\(\mathbb G_{a,k}\) acts trivially on it. In summary:

\paragraph{\textbf{Lemma 1.13.}}
\SGARefTarget{sga3:xx:1.13}
\emph{Under the hypotheses of \SGARefLink{sga3:xx:1.9}{1.9}, \(\mathfrak g\), as a
\(T\mathbin{\cdot_\alpha}\mathbb G_{a,k}\)-module, has a
composition series whose successive quotients are}
\[
  \mathfrak g^{-\alpha},\qquad \mathfrak t,\qquad
  \mathfrak g^\alpha,
\]
\emph{regarded as
\(T\mathbin{\cdot_\alpha}\mathbb G_{a,k}\)-modules through the
projection
\(T\mathbin{\cdot_\alpha}\mathbb G_{a,k}\to T\).}

We are thus reduced to computing the cohomology of
\(T\mathbin{\cdot_\alpha}\mathbb G_{a,k}\) acting on \(W(k)\)
through the projection
\(T\mathbin{\cdot_\alpha}\mathbb G_{a,k}\to T\) and a character
\(\beta\) of \(T\) (here \(\beta=0,\alpha\), or \(-\alpha\)).
\footnote{The editors added the sentence that follows.}
Write \(k[x_1,\ldots,x_n]\) for the polynomial algebra over \(k\)
in \(n\) variables, and \(k_q[x_1,\ldots,x_n]\) for its subspace
of homogeneous polynomials of degree \(q\).

\paragraph{\textbf{Lemma 1.14.}}
\SGARefTarget{sga3:xx:1.14}
\emph{With the preceding notation,}
\[
  \mathrm H^n\bigl(T\mathbin{\cdot_\alpha}\mathbb G_{a,k},k\bigr)
    =\mathrm H^n(C^\ast_{\alpha,\beta}),
\]
\emph{where the complex \(C^\ast_{\alpha,\beta}\) is defined by}
\[
  C^n_{\alpha,\beta}=
  \begin{cases}
    k_q[x_1,\ldots,x_n],
      &\text{if \(\beta=q\alpha\), with \(q\in\mathbb N^\ast\),}\\
    0,&\text{otherwise,}
  \end{cases}
\]
\emph{and}
\[
\begin{aligned}
(\delta f)(x_1,x_2,\ldots,x_{n+1})
  ={}&f(x_2,\ldots,x_{n+1})\\
   &+\sum_{i=1}^{n}(-1)^i
       f(x_1,\ldots,x_i+x_{i+1},\ldots,x_{n+1})\\
   &+(-1)^{n+1}f(x_1,\ldots,x_n).
\end{aligned}
\]

Indeed, by \SGARefLink{sga3:i:proposition:5-3-2}{Exposé I, 5.3.2}, the functor
\(M\mapsto\mathrm H^0(T,M)\) is exact on the category of
\(T\)-modules, and the \(\mathrm H^q(T,-)\) vanish. It follows, as
in the usual theory of group cohomology, that
\(\mathrm H^n(T\mathbin{\cdot_\alpha}\mathbb G_{a,k},k)\) may be
computed as the \(n\)-th cohomology group of the complex of
cochains of \(\mathbb G_{a,k}\) in \(k\) that are invariant under
\(T\), that is, those satisfying
\[
  f\bigl(\alpha(t)x_1,\ldots,\alpha(t)x_n\bigr)
    =\beta(t)f(x_1,\ldots,x_n).
\]
This is precisely the complex just displayed.

To prove \SGARefLink{sga3:xx:1.9}{1.9}, it is therefore enough to show that
\(\mathrm H^2(C^\ast_{\alpha,\beta})=0\) for
\(\beta=0,\alpha,-\alpha\), which is immediate.

\paragraph{\textbf{Remark 1.15.}}
\SGARefTarget{sga3:xx:1.15}
The groups \(\mathrm H^n(C^\ast_{\alpha,\beta})\) can be computed
explicitly when \(\beta=q\alpha\); see M.~Lazard,
\emph{Lois de groupes et analyseurs}, Annales E.N.S., 1955. In
particular,
\[
  \mathrm H^n(C^\ast_{\alpha,q\alpha})=0
  \qquad\text{for \(n>q\).}
\]

\paragraph{\textbf{Notations 1.16.}}
\SGARefTarget{sga3:xx:1.16}
The image of the canonical immersion
\[
  W(\mathfrak g^{-\alpha})\times_S T
    \times_S W(\mathfrak g^\alpha)\longrightarrow G
\]
will be denoted by \(\Omega\). It is an open subset of \(G\)
containing the identity section. The images of
\[
\begin{gathered}
  W(\mathfrak g^{-\alpha}),\qquad W(\mathfrak g^\alpha),\qquad
  W(\mathfrak g^{-\alpha})\times_S T,\qquad
  T\times_S W(\mathfrak g^\alpha)
\end{gathered}
\]
will be denoted, respectively,
\[
  U_{-\alpha},\qquad U_\alpha,\qquad
  U_{-\alpha}\mathbin{\cdot}T,\qquad T\mathbin{\cdot}U_\alpha.
\]
Then \(U_\alpha\) (respectively \(U_{-\alpha}\)) is a subgroup of
\(G\), canonically endowed with a vector-bundle structure, and
\[
  \intaut(t)(x)=x^{\alpha(t)}
  \qquad
  \bigl(\text{respectively }x^{-\alpha(t)}\bigr)
\]
for every \(S'\to S\), \(t\in T(S')\), and
\(x\in U_\alpha(S')\) (respectively \(x\in U_{-\alpha}(S')\)).

There are canonical isomorphisms
\[
  T\mathbin{\cdot}U_\alpha\simeq T\mathbin{\cdot_\alpha}U_\alpha,
  \qquad
  T\mathbin{\cdot}U_{-\alpha}
    \simeq T\mathbin{\cdot_{-\alpha}}U_{-\alpha}.
\]
The open subset \(\Omega\) is stable under \(\intaut(T)\), and
\[
  \intaut(t')\bigl(y\mathbin{\cdot}t\mathbin{\cdot}x\bigr)
    =y^{-\alpha(t')}\mathbin{\cdot}t
       \mathbin{\cdot}x^{\alpha(t')}.
\]

\paragraph{\textbf{Corollary 1.17.}}
\SGARefTarget{sga3:xx:1.17}
\emph{We have}
\[
  \Lie(U_\alpha/S)=\mathfrak g^\alpha,\qquad
  \Lie(U_{-\alpha}/S)=\mathfrak g^{-\alpha}.
\]
\emph{The isomorphisms}
\[
  W(\mathfrak g^\alpha)\xrightarrow[\sim]{\exp}U_\alpha,
  \qquad
  W(\mathfrak g^{-\alpha})\xrightarrow[\sim]{\exp}U_{-\alpha}
\]
\emph{are those of \SGARefLink{sga3:xix:lemma:4:2}{Exposé XIX, 4.2}.}

\paragraph{\textbf{Corollary 1.18.}}
\SGARefTarget{sga3:xx:1.18}
\emph{The open subset \(\Omega\) is relatively schematically dense
in \(G\) \textup{(cf. Exposé XVIII, \S1)}.}

This follows immediately from Exposé XVIII, 1.3.

\paragraph{\textbf{Corollary 1.19.}}
\SGARefTarget{sga3:xx:1.19}
\emph{The center of \(G\) is}
\[
  \operatorname{Centr}(G)=\Ker(\alpha).
\]
\emph{It is therefore a closed subgroup of \(G\), of
multiplicative type and of finite type.}

The second assertion follows from the first by \SGARefLink{sga3:ix:proposition:2:7}{Exposé IX, 2.7}. Let
us prove the first. The inner automorphism defined by a section of
\(\Ker(\alpha)\) acts trivially on \(\Omega\), by the last formula
of \SGARefLink{sga3:xx:1.16}{1.16}, and hence on \(G\), by \SGARefLink{sga3:xx:1.18}{1.18}. Conversely, if
\(g\in G(S)\) centralizes \(G\), then it centralizes \(T\) and
\(U_\alpha\). It is therefore a section of \(T\), by \SGARefLink{sga3:xix:paragraph:2:8}{Exposé XIX,
2.8}, and annihilates \(\alpha\). Since the same argument applies
after every base change, we indeed have
\(\operatorname{Centr}(G)=\Ker(\alpha)\).

\paragraph{\textbf{Corollary 1.20.}}
\SGARefTarget{sga3:xx:1.20}
\emph{A monomorphism
\[
  p_\alpha\colon\mathbb G_{a,S}\longrightarrow G
\]
normalized by \(T\), with multiplier \(\alpha\), exists if and
only if the \(\mathscr O_S\)-module \(\mathfrak g^\alpha\) is
free. More precisely, the assignments}
\[
  X_\alpha\longmapsto
    \bigl(x\longmapsto\exp(xX_\alpha)\bigr),
  \qquad
  p_\alpha\longmapsto\Lie(p_\alpha)
\]
\emph{give a bijection between
\(\Gamma(S,\mathfrak g^\alpha)^\times\) and the set of such
monomorphisms \(p_\alpha\), which is also the set of isomorphisms
of vector group schemes}
\[
  \mathbb G_{a,S}\xrightarrow{\sim}U_\alpha.
\]

\paragraph{\textbf{Corollary 1.21.}}
\SGARefTarget{sga3:xx:1.21}
\emph{The subgroups \(U_\alpha\) and \(U_{-\alpha}\) of \(G\)
commute on no fiber.}

Indeed, if \((U_\alpha)_s\) and \((U_{-\alpha})_s\) commute, then
\(\Omega_s\) is a subgroup of \(G_s\), and hence
\(\Omega_s=G_s\). It follows that \(G_s\) is
solvable, contrary to the hypothesis that \(G_s\) is reductive of
semisimple rank \(1\).
